\section*{Chapitre \ref{ch:dv:asians-lookbacks-cliquets-and-forward-starts} --- Options asiatiques, lookbacks, cliquets et options à départ différé}

\begin{solution}{exo:dv:asians-lookbacks-cliquets-and-forward-starts:1}
$n=4$~: $\sqrt{5\times9/(6\times16)}=\sqrt{0.469}=0.685$. Lorsque $n\to\infty$, le facteur tend vers $\sqrt{2n^2/6n^2}=1/\sqrt3=0.577$.
\end{solution}

\begin{solution}{exo:dv:asians-lookbacks-cliquets-and-forward-starts:2}
Trajectoire par trajectoire, la moyenne arithmétique de nombres positifs est au moins égale à leur moyenne géométrique, et le payoff de l'option d'achat croît
avec la moyenne~; le payoff arithmétique domine donc le payoff géométrique sur chaque trajectoire.
\end{solution}

\begin{solution}{exo:dv:asians-lookbacks-cliquets-and-forward-starts:3}
Aujourd'hui, l'option d'achat à six mois à la monnaie vaut $S\,C_{\mathrm{BS}}(1,1,0.5)$ par homogénéité~; l'\omterm{def:dv:asians-lookbacks-cliquets-and-forward-starts:fwdstart}{option à départ différé} vaut
$Se^{-q\times0.5}C_{\mathrm{BS}}(1,1,0.5)$. Le rapport est $e^{-0.01}=0.990$~: les dividendes versés avant la fixation du prix d'exercice abaissent
le spot auquel il est fixé.
\end{solution}

\begin{solution}{exo:dv:asians-lookbacks-cliquets-and-forward-starts:4}
Si $r\le f$~: $f+0-0=f$. Si $f<r<c$~: $f+(r-f)-0=r$. Si $r\ge c$~: $f+(r-f)-(r-c)=c$. Le rendement écrêté est donc une constante plus un
call spread entre $f$ et $c$, $(r-f)^+-(r-c)^+$. Lorsque $c-f$ est petit, ce call spread vaut $(c-f)$ fois un digital à environ $(f+c)/2$~: chaque période
rapporte comme un digital à départ différé, dont le prix est fixé par le skew forward.
\end{solution}

\begin{solution}{exo:dv:asians-lookbacks-cliquets-and-forward-starts:5}
La \omterm{def:dv:local-volatility:model}{volatilité locale} ajuste le skew d'aujourd'hui au moyen d'une volatilité qui est une fonction fixe du spot et du temps~; à mesure que la date forward
s'éloigne, le spot a bougé et le skew local rencontré est celui calibré pour décroître avec la maturité, si bien que le
\omterm{def:dv:local-volatility:forward}{smile forward} s'aplatit. Dans le \omterm{def:dv:stochastic-volatility:heston}{modèle de Heston}, le skew provient de la corrélation du spot avec une variance encore aléatoire à la
date forward~: elle y est régénérée, et le caractère aléatoire de la variance ajoute de la courbure.
\end{solution}

\begin{solution}{exo:dv:asians-lookbacks-cliquets-and-forward-starts:6}
L'écart type de l'estimateur brut par trajectoire vaut environ $0.025\sqrt{100\,000}=8.0$, si bien qu'une erreur de 0.001 exige
$(8.0/0.001)^2=64$ millions de trajectoires~; avec la variable de contrôle, environ 47\,000. Le calcul du payoff géométrique ajoute peu à chaque trajectoire,
si bien que le temps de calcul baisse à peu près dans le même rapport.
\end{solution}

\begin{solution}{exo:dv:asians-lookbacks-cliquets-and-forward-starts:7}
Avec des coupons de 15~\%, 25~\% et 35~\%~: \omterm{def:dv:local-volatility:model}{volatilité locale} 1.23~\%, 5.43~\% et 12.06~\%~; Heston 2.49~\%, 7.56~\% et 14.55~\%~; \omterm{def:dv:black-scholes-three-ways:model}{modèle de Black--Scholes}
0.58~\%, 3.57~\% et 10.00~\%. Le prix de Heston varie le plus en points (12.1 contre 10.8 et 9.4)~; au coupon faible, il vaut deux fois le
prix du \omterm{def:dv:local-volatility:model}{modèle de volatilité locale}, car il accorde davantage de poids aux trajectoires comportant beaucoup de fortes baisses.
\end{solution}

\begin{solution}{exo:dv:asians-lookbacks-cliquets-and-forward-starts:8}
Les options vanilles fixent les distributions terminales à chaque échéance, et non la loi conjointe des rendements entre elles. Un \omterm{def:dv:asians-lookbacks-cliquets-and-forward-starts:cliquet}{cliquet} paie sur
rendements forward, si bien que son prix dépend du \omterm{def:dv:local-volatility:forward}{smile forward}, que le \omterm{def:dv:local-volatility:model}{modèle de volatilité locale} prévoit à partir de ses propres
hypothèses. Un autre modèle calibré sur les mêmes vanilles valorise le \omterm{def:dv:asians-lookbacks-cliquets-and-forward-starts:cliquet}{cliquet} du chapitre 45~\% plus cher~; l'ajustement ne constitue pas une preuve en faveur de
l'un ou l'autre prix.
\end{solution}

\begin{solution}{pb:dv:asians-lookbacks-cliquets-and-forward-starts:1}
\textbf{1.} 5.32 et 8.83.
\textbf{2.} 0.376~; $0.61\times20\%=12.3\%$.
\textbf{3.} 0.025 sans et 0.0007 avec la variable de contrôle.
\textbf{4.} 15.69 en continu~; en quotidien, 15.09 par simulation et 15.08 par la formule décalée.
\textbf{5.} La correction est un développement en $\sigma\sqrt{\Delta t}$~; avec un trimestre entre les dates, cela donne 0.1, trop grand, et
le minimum entre les dates est loin de celui relevé aux dates de constatation.
\textbf{6.} \omterm{def:dv:local-volatility:model}{Volatilité locale} 19.19~\%, Heston 18.99~\%, marché 19.19~\%.
\textbf{7.} \omterm{def:dv:local-volatility:model}{Volatilité locale} 21.5~\%, Heston 18.5~\%~; la volatilité à un mois d'aujourd'hui 14.9~\%.
\textbf{8.} \omterm{def:dv:local-volatility:model}{Volatilité locale} $23.9-21.3=2.6$ points, Heston $24.0-20.0=4.0$~; la volatilité à un mois d'aujourd'hui $19.6-12.0=7.6$.
\textbf{9.} La structure par terme de la surface est croissante~: la variance entre six et sept mois est plus élevée que sur le premier
mois.
\textbf{10.} Celui dont le skew forward correspond aux cotations des \omterm{def:dv:asians-lookbacks-cliquets-and-forward-starts:fwdstart}{options à départ différé}~; un \omterm{def:dv:stochastic-volatility:sv}{modèle à volatilité stochastique} régénère le skew,
ce que ne fait pas la \omterm{def:dv:local-volatility:model}{volatilité locale}, bien que le \omterm{def:dv:rough-volatility-and-forward-variance-models:bergomi}{modèle de Bergomi} ait constaté que Heston surpondère probablement les scénarios
à faible volatilité et à fort skew pour les \omterm{def:dv:asians-lookbacks-cliquets-and-forward-starts:fwdstart}{options à départ différé} courtes.
\textbf{11.} 1.53~\%, 2.22~\% et 1.19~\% du notionnel.
\textbf{12.} Pour $\pm0.5\%$~: 0.78~\%, 1.16~\% et 0.61~\%. Pour $\pm5\%$~: 5.73~\%, 6.12~\% et 4.74~\%.
\textbf{13.} Un cap large rend le payoff de chaque période presque égal au rendement lui-même, linéaire, dont la valeur ne dépend pas du
skew forward~; le prix reflète alors le niveau de la volatilité forward et le floor global, sur lequel les modèles divergent moins.
\textbf{14.} 5.43~\%, 7.56~\% et 3.57~\%.
\textbf{15.} Le vendeur est short de puts forward et préférerait le prix moins cher en \omterm{def:dv:local-volatility:model}{volatilité locale}~; c'est là l'alerte,
car le modèle qui rend le produit bon marché est celui dont le \omterm{def:dv:local-volatility:forward}{smile forward} a perdu son skew.
\textbf{16.} Les \omterm{def:dv:asians-lookbacks-cliquets-and-forward-starts:fwdstart}{options à départ différé} et les \omterm{def:dv:asians-lookbacks-cliquets-and-forward-starts:cliquet}{cliquets} eux-mêmes lorsqu'ils sont cotés~; les options sur l'indice de volatilité, qui révèlent la
distribution de la volatilité forward~; la \omterm{def:dv:rough-volatility-and-forward-variance-models:fwdvar}{variance forward} issue des \omterm{def:dv:variance-swaps-and-volatility-derivatives:vs}{swaps de variance}.
\textbf{17.} Avec des risk reversals à départ différé là où ils se négocient~; sinon avec des risk reversals vanilles de deux échéances,
dimensionnés par le modèle, et une réserve pour l'erreur du modèle.
\textbf{18.} Au moins une large part de l'écart avec le modèle alternatif~: 0.69 point de notionnel sur cette transaction.
\textbf{19.} 0.69 point de notionnel~: 2.22~\% selon Heston contre 1.53~\% selon la \omterm{def:dv:local-volatility:model}{volatilité locale}, soit 45~\% de plus, les deux modèles étant calibrés sur
la même surface.
\textbf{20.} Le \omterm{def:dv:local-volatility:forward}{smile forward}~: la forme que prendra le smile à des dates futures, qu'aucune option vanille ne rémunère.
\end{solution}

\begin{solution}{iq:dv:asians-lookbacks-cliquets-and-forward-starts:1}
La moyenne fluctue moins que le prix final~: sa variance vaut environ le tiers de la variance terminale dans le cas d'une moyenne continue
$(n+1)(2n+1)/6n^2$ pour $n$ relevés, si bien qu'une \omterm{def:dv:asians-lookbacks-cliquets-and-forward-starts:asian}{option asiatique} à la monnaie vaut à peu près $1/\sqrt3\approx58\%$ de l'option vanille
(60~\% avec douze relevés mensuels).
\iqlookfor{la réduction de variance et le facteur.}
\end{solution}

\begin{solution}{iq:dv:asians-lookbacks-cliquets-and-forward-starts:2}
Simuler des trajectoires aux dates de constatation, calculer les payoffs arithmétique et géométrique sur les mêmes trajectoires, et utiliser la
formule fermée de l'\omterm{def:dv:asians-lookbacks-cliquets-and-forward-starts:asian}{option asiatique} géométrique comme variable de contrôle, avec un coefficient de régression~; ajouter des tirages antithétiques~; rapporter
l'erreur type. La réduction de variance est typiquement d'un facteur mille.
\iqlookfor{la variable de contrôle et le rapport d'erreur.}
\end{solution}

\begin{solution}{iq:dv:asians-lookbacks-cliquets-and-forward-starts:3}
Par homogénéité~: en $t_1$ elle vaut $S_{t_1}C_{\mathrm{BS}}(1,k,t_2-t_1)$, donc aujourd'hui $S_0e^{-qt_1}C_{\mathrm{BS}}(1,k,t_2-t_1)$. Dans un
modèle à smile, elle dépend du \omterm{def:dv:local-volatility:forward}{smile forward}, soit la prévision faite par le modèle pour le smile en $t_1$ sur la période jusqu'à $t_2$.
\iqlookfor{homogénéité, et le \omterm{def:dv:local-volatility:forward}{smile forward}.}
\end{solution}

\begin{solution}{iq:dv:asians-lookbacks-cliquets-and-forward-starts:4}
Au \omterm{def:dv:local-volatility:forward}{smile forward}~: chaque période est un call spread à départ différé moins un put spread, proche d'un digital lorsque les caps sont
serrés, si bien que le skew forward en fixe la valeur. Deux modèles peuvent reproduire toutes les options vanilles (distributions terminales) et pourtant attribuer
des lois conjointes différentes aux rendements forward.
\iqlookfor{skew forward, et distributions terminales contre distributions conjointes.}
\end{solution}

\begin{solution}{iq:dv:asians-lookbacks-cliquets-and-forward-starts:5}
Son skew est une fonction fixe du spot et du temps calibrée aujourd'hui, et ses \omterm{def:dv:local-volatility:forward}{smiles forward} s'aplatissent~: il sous-évalue le skew forward
et la \omterm{def:dv:stochastic-volatility:sv}{volatilité de la volatilité}, si bien qu'il sous-évalue les \omterm{def:dv:asians-lookbacks-cliquets-and-forward-starts:cliquet}{cliquets} et les \omterm{def:dv:asians-lookbacks-cliquets-and-forward-starts:reverse}{reverse cliquets} (de 45~\% par rapport à Heston dans
l'exemple du chapitre).
\iqlookfor{aplatissement des \omterm{def:dv:local-volatility:forward}{smiles forward} et sens de l'erreur.}
\end{solution}

\begin{solution}{iq:dv:asians-lookbacks-cliquets-and-forward-starts:6}
Positions courtes en puts forward sur de nombreuses périodes~: exposition à la volatilité forward, au skew forward et à la \omterm{def:dv:stochastic-volatility:sv}{volatilité de la
volatilité}, et aux krachs qui frappent plusieurs périodes à la fois. Couvrir le \omterm{def:dv:greeks-and-the-hedging-pnl:greeks}{vega} par compartiment forward au moyen de \omterm{def:dv:variance-swaps-and-volatility-derivatives:vs}{swaps de variance} ou de produits à départ différé,
le skew par des risk reversals, et conserver des réserves pour l'écart de modèle et pour la \omterm{def:dv:stochastic-volatility:sv}{volatilité de la volatilité}, qui ne dispose
d'aucune couverture liquide.
\iqlookfor{les expositions, les couvertures par compartiment et la part non couvrable.}
\end{solution}
